\section{Calcul et sélection des actionneurs}
En raison de la symétrie de l’architecture mécanique, le dimensionnement
des actionneurs peut être réduit à deux cas distincts. Les moteurs $M_1$ et $M_3$
assurent la même fonction sur deux sous-ensembles symétriques et sont donc soumis
aux mêmes efforts. De même, les moteurs $M_2$ et $M_4$ entraînent les seconds bras
des deux mécanismes et peuvent être dimensionnés de manière identique.
%ajouter une figure avec les moteurs lequels sont les moteurs 1,2,3,4
\begin{figure}[H]
\centering
\includegraphics[width=0.8\textwidth]{images/robot_scara.png}
\caption{Identification des moteurs $M_1$, $M_2$, $M_3$ et $M_4$}
\label{fig:moteurs}
\end{figure}
Ainsi, deux dimensionnements sont nécessaires :
\begin{itemize}
    \item le dimensionnement du moteur de premier axe, appliqué à $M_1$ et $M_3$ ;
    \item le dimensionnement du moteur de second axe, appliqué à $M_2$ et $M_4$.
\end{itemize}
\subsection{Hypothèses de dimensionnement}

Le robot évolue principalement dans un plan horizontal. Les axes entraînés par
les moteurs $M_1$, $M_2$, $M_3$ et $M_4$ sont des axes de rotation verticaux.
Dans ces conditions, l’influence du poids sur le couple moteur autour de ces axes
peut être négligée. Le dimensionnement est donc principalement gouverné par :
\begin{itemize}
    \item l’inertie des bras et des éléments embarqués ;
    \item l’accélération angulaire imposée ;
    \item les pertes dans la transmission ;
    \item une marge de sécurité permettant de garantir un fonctionnement fiable.
\end{itemize}

Le couple nécessaire est évalué selon la relation suivante :
\begin{equation}
    C_{\mathrm{req}} = C_{\mathrm{dyn}} + C_{\mathrm{res}}
\end{equation}
où :
\begin{itemize}
    \item $C_{\mathrm{dyn}}$ représente le couple dynamique nécessaire à l’accélération ;
    \item $C_{\mathrm{res}}$ représente un terme additionnel regroupant les frottements
    et les incertitudes de modélisation.
\end{itemize}

Le couple dynamique s’exprime par :
\begin{equation}
    C_{\mathrm{dyn}} = J_{\mathrm{eq}} \, \alpha
\end{equation}
avec :
\begin{itemize}
    \item $J_{\mathrm{eq}}$ le moment d’inertie équivalent ramené à l’axe moteur ;
    \item $\alpha$ l’accélération angulaire maximale.
\end{itemize}

En présence d’une transmission de rapport $i$ et de rendement $\eta$, le couple
à fournir par le moteur devient :
\begin{equation}
    C_{\mathrm{mot}} = \frac{C_{\mathrm{req}}}{i \, \eta}
\end{equation}

Enfin, le moteur retenu doit vérifier :
\begin{equation}
    C_{\mathrm{nom}} \geq S \, C_{\mathrm{mot}}
\end{equation}
où $S$ désigne un coefficient de sécurité, généralement compris entre $1.5$ et $2$.
\subsection{Moteurs des axes principaux}
Les moteurs $M_1$ et $M_3$ entraînent les premiers bras des deux mécanismes
symétriques. Ils doivent donc fournir le couple nécessaire pour accélérer
l’ensemble mobile situé en aval du premier axe.

Le couple dynamique est donné par :
\begin{equation}
    C_{\mathrm{dyn},13} = J_{\mathrm{eq},13} \cdot \alpha_{13}
\end{equation}

où $J_{\mathrm{eq},13}$ est le moment d’inertie équivalent autour de l’axe
principal et $\alpha_{13}$ l’accélération angulaire maximale.

En tenant compte du rapport de transmission $i_{13}$, du rendement $\eta_{13}$
et d’un coefficient de sécurité $S$, le couple minimal du moteur est :
\begin{equation}
    C_{\mathrm{mot},13} =
    S \cdot \frac{J_{\mathrm{eq},13} \cdot \alpha_{13}}{i_{13}\eta_{13}}
\end{equation}

Par symétrie :
\begin{equation}
    C_{M_1} = C_{M_3} = C_{\mathrm{mot},13}
\end{equation}
\subsection{Moteurs des axes secondaires}

Les moteurs $M_2$ et $M_4$ entraînent les seconds bras. Ils sont soumis à une
inertie plus faible, car seuls les éléments situés après le second axe sont pris
en compte.

Le couple dynamique est :
\begin{equation}
    C_{\mathrm{dyn},24} = J_{\mathrm{eq},24} \cdot \alpha_{24}
\end{equation}

Le couple minimal requis au niveau moteur est donc :
\begin{equation}
    C_{\mathrm{mot},24} =
    S \cdot \frac{J_{\mathrm{eq},24} \cdot \alpha_{24}}{i_{24}\eta_{24}}
\end{equation}

Par symétrie :
\begin{equation}
    C_{M_2} = C_{M_4} = C_{\mathrm{mot},24}
\end{equation}
\subsection{Sélection des moteurs}

Les moteurs sont sélectionnés en vérifiant que leur couple nominal est supérieur
au couple minimal calculé :
\begin{equation}
    C_{\mathrm{nom}} \geq C_{\mathrm{mot}}
\end{equation}

La vitesse maximale du moteur doit également être compatible avec la vitesse
angulaire souhaitée en sortie de transmission :
\begin{equation}
    \omega_{\mathrm{mot}} = i \cdot \omega_{\mathrm{sortie}}
\end{equation}

Le choix final est résumé dans le tableau~\ref{tab:selection_moteurs}.
